\clearpage
\appendix
\chapter*{Annexes}
\addcontentsline{toc}{chapter}{Annexes}
\markboth{Annexes}{Annexes}
\renewcommand{\labelitemi}{$\bullet$}

\section*{Annexe A : Stack de Conteneurisation (\texttt{docker-compose.yml})}
\addcontentsline{toc}{section}{Annexe A : Stack de Conteneurisation (docker-compose.yml)}
Le fichier \texttt{docker-compose.yml} ci-dessous orchestre l'ensemble des 5 microservices conteneurisés de la plateforme au sein du réseau privé virtuel \texttt{rag\_network} :

\begin{verbatim}
version: '3.8'

services:
  qdrant_db:
    image: qdrant/qdrant:latest
    container_name: qdrant_db
    ports:
      - "6333:6333"
      - "6334:6334"
    volumes:
      - qdrant_data:/qdrant/storage
    networks:
      - rag_network

  ollama-llm:
    image: ollama/ollama:latest
    container_name: ollama-llm
    ports:
      - "11434:11434"
    volumes:
      - ollama_data:/root/.ollama
    networks:
      - rag_network

  chatbot_api:
    build: .
    container_name: chatbot_api
    ports:
      - "8000:8000"
    environment:
      - QDRANT_HOST=qdrant_db
      - QDRANT_PORT=6333
      - OLLAMA_BASE_URL=http://ollama-llm:11434
    volumes:
      - ./data:/app/data
    networks:
      - rag_network
    depends_on:
      - qdrant_db
      - ollama-llm

  mlflow_server:
    image: ghcr.io/mlflow/mlflow:v2.10.0
    container_name: mlflow_server
    ports:
      - "5001:5000"
    networks:
      - rag_network

  langfuse_server:
    image: langfuse/langfuse:2
    container_name: langfuse_server
    ports:
      - "3000:3000"
    networks:
      - rag_network

networks:
  rag_network:
    driver: bridge

volumes:
  qdrant_data:
  ollama_data:
\end{verbatim}

\clearpage

\section*{Annexe B : Structure du Dataset de Benchmark RAG (\texttt{eval\_dataset.json})}
\addcontentsline{toc}{section}{Annexe B : Structure du Dataset de Benchmark RAG (eval\_dataset.json)}
Le dataset de test \texttt{tests/eval\_dataset.json} (versionné avec DVC sous \texttt{eval\_dataset.json.dvc}) contient 50 cas d'évaluation. L'extrait ci-dessous illustre la structure des questions intra-domaine et des questions pièges hors-domaine :

\begin{verbatim}
[
  {
    "id": "q01_intradomaine",
    "type": "intra_domain",
    "question": "Quelle est la politique de télétravail chez Owl Media ?",
    "ground_truth": "Le télétravail est autorisé à raison de 2 jours par semaine après validation du responsable hiérarchique."
  },
  {
    "id": "q15_piege_horsdomaine",
    "type": "out_of_domain",
    "question": "Quel est le chiffre d'affaires mondial de Microsoft en 2020 ?",
    "ground_truth": "NO_CONTEXT"
  }
]
\end{verbatim}

\clearpage

\section*{Annexe C : Schéma de la Base de Données de Supervision (\texttt{monitoring.db})}
\addcontentsline{toc}{section}{Annexe C : Schéma de la Base de Données de Supervision (monitoring.db)}
La table SQLite \texttt{queries} dans la base locale \texttt{data/monitoring.db} consigne de manière asynchrone l'ensemble des métriques de latence et d'exécution décomposées :

\begin{verbatim}
CREATE TABLE IF NOT EXISTS queries (
    query_id TEXT PRIMARY KEY,
    timestamp DATETIME DEFAULT CURRENT_TIMESTAMP,
    question TEXT NOT NULL,
    retrieval_mode TEXT NOT NULL,
    top_k INTEGER NOT NULL,
    reranking_enabled BOOLEAN NOT NULL,
    retrieval_latency REAL NOT NULL,
    hybrid_latency REAL NOT NULL,
    reranking_latency REAL NOT NULL,
    llm_latency REAL NOT NULL,
    total_latency REAL NOT NULL,
    candidates_count INTEGER NOT NULL,
    final_sources_count INTEGER NOT NULL,
    status TEXT NOT NULL,
    error_type TEXT
);
\end{verbatim}

\clearpage

\section*{Annexe D : Galerie et Description des Interfaces Applicatives}
\addcontentsline{toc}{section}{Annexe D : Galerie et Description des Interfaces Applicatives}
Cette annexe présente le récapitulatif visuel et descriptif des différentes vues de la plateforme Owl Media RAG MLOps :

\begin{enumerate}
    \item \textbf{Interface Conversationnelle (Chatbot SPA)} : Interface Web principale d'interaction en langage naturel. Elle permet la saisie fluide des requêtes utilisateur, la réception de la réponse générée et l'inspection des extraits sources certifiés avec leur score de pertinence.
    \item \textbf{Panneau d'Administration Documentaire} : Console dédiée au rôle \texttt{Admin} offrant le chargement par glisser-déposer des documents (PDF, DOCX, TXT), le déclenchement du pipeline d'ingestion/vectorisation et la suppression de fichiers de la base vectorielle Qdrant.
    \item \textbf{Tableau de Bord de Supervision MLOps (SQLite)} : Interface de suivi analytique affichant les percentiles de latence (P50/P95), l'historique chronologique des requêtes et le taux de blocage anti-hallucination.
    \item \textbf{Console d'Observabilité Distribuée (Langfuse)} : Interface d'inspection visuelle détaillée montrant l'arbre d'exécution des prompts et des sous-agents rattachés à chaque \texttt{query\_id}.
    \item \textbf{Plateforme de Tracking d'Expérimentations (MLflow)} : Dashboard de suivi des campagnes d'évaluation automatisées comparant les métriques RAG (Fidélité, Hit Rate, Pertinence) entre différentes configurations de recherche.
\end{enumerate}
